\section{Proofs}
\label{app:proofs}

\subsection{Proof of Theorem~\ref{thm:validity}}
Fix the candidate order $\pi_{(1)}, \pi_{(2)}, \dots$ (a function of
the train split only, hence independent of the calibration sample).
For candidate $\pi$ let $H_0(\pi) = \bigcup_j \{R_j(\pi) > \alpha_j\}$
be its composite null. The per-candidate test rejects iff every
constraint's Hoeffding--Bentkus p-value is $\le \delta$. Each
$p^{(j)}_\pi$ is a valid p-value for $H_0^{(j)}$
\citep{bentkus2004hoeffding,angelopoulos2021ltt}; by the
intersection--union principle \citep{berger1982iut}, when $H_0(\pi)$
holds at least one $H_0^{(j)}$ holds, so
$\Prob(\text{reject } \pi) \le \Prob(p^{(j^\ast)}_\pi \le \delta) \le
\delta$ for a true null $j^\ast$: the composite test has size
$\le\delta$. Let $k^\ast$ be the index of the first infeasible
candidate in the order (if none, no error is possible). The procedure
deploys an infeasible candidate only if the walk certifies
$\pi_{(k^\ast)}$ (it cannot reach later candidates without doing so,
and all earlier candidates are feasible). That event has probability
$\le \delta$. Choosing the cheapest \emph{certified} candidate only
ever selects among candidates that were certified, so the error event
is unchanged. \hfill$\square$

\subsection{Proof of Proposition~\ref{prop:composition}}
Write $\rho_i^{\mathrm{step}} = \Prob_q(\ell(P_i, q) >
\ell(P_{i-1}, q))$ for the stepwise harm of $R_i$ against its
predecessor. \emph{Deterministic inclusion}: for any query $q$ with
$\ell(P_m, q) > \ell(P_0, q)$, the telescoping sum
$\sum_{i=1}^m [\ell(P_i, q) - \ell(P_{i-1}, q)] =
\ell(P_m, q) - \ell(P_0, q) > 0$ forces at least one strictly positive
summand, i.e., $q \in \{\ell(P_i) > \ell(P_{i-1})\}$ for some $i$.
Hence $\{\ell(P_m) > \ell(P_0)\} \subseteq \bigcup_i \{\ell(P_i) >
\ell(P_{i-1})\}$ pointwise, with no assumption on the rewrites. Taking
probabilities,
$\rho(R_m \cdots R_1; P_0) \le \sum_i \rho_i^{\mathrm{step}}$. Each
step voucher satisfies $\Prob(\rho_i^{\mathrm{step}} >
\hat\eps_i^{\mathrm{step}}) \le \delta_r$ (Clopper--Pearson); on the
event that all $m$ bounds hold (probability $\ge 1 - m\delta_r$ by a
union bound over the $m$ vouchers, the only stochastic step),
$\rho(R_m \cdots R_1; P_0) \le \sum_i \hat\eps_i^{\mathrm{step}}$.
\hfill$\square$

\begin{rmk}
The base-relative union bound (each rewrite vouchered against $P_0$)
requires the additional no-synergistic-harm assumption
$\{\ell(P_m) > \ell(P_0)\} \subseteq \bigcup_i \{\ell(R_i(P_0)) >
\ell(P_0)\}$, which composed rewrites violate on real workloads
(Section~\ref{sec:exp-rewrites}); the telescoping decomposition
replaces that behavioral assumption with an algebraic identity by
conditioning each voucher on the actual chain prefix.
\end{rmk}

\subsection{Proof of Proposition~\ref{prop:eprocess}}
Under the null $\E[\ell_t \mid \mathcal{F}_{t-1}] \le \alpha$, each
per-step factor satisfies
$\E[1 + w(\ell_t - \alpha) \mid \mathcal{F}_{t-1}]
= 1 + w(\E[\ell_t \mid \mathcal{F}_{t-1}] - \alpha) \le 1$
for any $w \in (0, 1/\alpha)$ and $\ell_t \in [0,1]$ (the factor is
also nonnegative since $w < 1/\alpha$ and $\ell_t \ge 0$). Hence each
restarted product
$E_t^{(r,w)} = \prod_{i=r}^{t}(1 + w(\ell_i - \alpha))$ is a
nonnegative supermartingale, equal to $1$ for $t < r$, so
$\E[E_t^{(r,w)} \mid \mathcal{F}_{t-1}] \le E_{t-1}^{(r,w)}$ for all
$t$. The statistic in~\eqref{eq:mixdetector} is a fixed convex
combination of these supermartingales and the constants $1$ (the
not-yet-started restarts $r > t$, each contributing
$\tfrac{1}{T|W|}\cdot 1$, which total $1 - t/T$). A convex combination
of nonnegative supermartingales, with weights fixed in advance and
summing to one, is again a nonnegative supermartingale, and
$M_0 = 1$. This is exactly where the mixture is needed: a
\emph{maximum} $\max_{r} E_t^{(r,w)}$ is not a supermartingale (the
maximizing index is chosen with hindsight), so Ville's inequality
would not apply to it. Applying Ville's inequality
\citep{ville1939etude} to $M_t$ gives
$\Prob(\exists t \le T : M_t \ge 1/\delta)
= \Prob(\sup_{t \le T} M_t \ge 1/\delta) \le \E[M_0]\,\delta =
\delta$. \hfill$\square$

\subsection{Proof of Proposition~\ref{prop:delay}}
Consider the single mixture component with restart $r = \tau + 1$ and
the growth-optimal grid bet $w^\ast$, the maximizer of
$\E[\log(1 + w(\ell_t - \alpha))]$ over $w \in W$,
whose post-change per-step
log-growth is $g^\ast > 0$. Its increments
$X_i = \log\big(1 + w^\ast(\ell_i - \alpha)\big)$ are bounded
($|X_i| \le \log\frac{1}{1-w^\ast\alpha} \vee
\log(1 + w^\ast(1-\alpha)) =: b$) with conditional mean $\ge g^\ast$
for $i > \tau$. By the Azuma--Hoeffding bound applied to
$\sum_{i=\tau+1}^{\tau+n}(X_i - g^\ast)$, with probability
$\ge 1 - \delta'$,
\[
\log E_{\tau+n}^{(\tau+1, w^\ast)}
= \sum_{i=\tau+1}^{\tau+n} X_i
\ \ge\ n g^\ast - b\sqrt{2 n \log(1/\delta')}.
\]
Since $M_{\tau+n} \ge \frac{1}{T|W|}E_{\tau+n}^{(\tau+1,w^\ast)}$, the
alarm threshold $M_{\tau+n} \ge 1/\delta$ is reached once
$n g^\ast - b\sqrt{2 n \log(1/\delta')} \ge \log(1/\delta) +
\log(T|W|)$. Solving this scalar inequality gives
\[
n = O\big((\log(1/\delta) + \log(T|W|) +
\log(1/\delta'))/g^\ast\big),
\]
independent of $\tau$. For a Bernoulli
loss the Kelly bet
$w^\ast = \frac{\alpha'-\alpha}{\alpha'(1-\alpha)}$ attains
$g^\ast = \mathrm{KL}(\alpha' \Vert \alpha)$; any grid $W$ bracketing
it loses only a constant factor, giving the stated
$\mathrm{KL}(\alpha'\Vert\alpha)$ rate. \hfill$\square$

\subsection{Proof of Proposition~\ref{prop:consistency}}
Let $\pi^\dagger$ be the cheapest candidate with
$R(\pi^\dagger) \le \alpha - \gamma$, at position $K$ in the walk,
with all true risks non-decreasing along the order up to $K$.
\emph{Certification of the prefix.} For every $\pi_{(k)}$, $k \le K$:
$R(\pi_{(k)}) \le \alpha - \gamma$, and by Hoeffding the empirical
risk is $\le \alpha - \gamma/2$ with probability
$1 - e^{-n\gamma^2/2}$; the Hoeffding--Bentkus p-value at empirical
risk $\alpha - \gamma/2$ is
$\le e^{-n\,h_1(\alpha-\gamma/2,\alpha)} \le e^{-n\gamma^2/2}$, which
is $\le \delta$ once $n \ge 2\gamma^{-2}\log(1/\delta)$. A union bound
over the $K$ candidates (at level $\delta'/K$ each) shows all are
certified simultaneously with probability $\ge 1 - \delta'$ when
$n \ge 2\gamma^{-2}(\log(1/\delta) \vee \log(K/\delta'))$; on this
event the walk reaches at least position $K$ and $\pi^\dagger$ is in
the certified set. \emph{Uniform cost estimation.} The walk does not
stop at $K$ --- it continues until the first failed test --- so the
deployed policy may sit at any position in the certified set, whose
extent is data-dependent and not known in advance. We therefore
control cost estimates uniformly over \emph{all} $M$ candidates:
per-query costs lie in $[0, C_{\max}]$, so by Hoeffding plus a union
bound over the $M$ candidates (at level $\delta'/M$ each), with
probability $\ge 1 - \delta'$ every candidate's estimated mean
calibration cost $\hat c(\pi)$ satisfies
$|\hat c(\pi) - c(\pi)| \le \eps_n :=
C_{\max}\sqrt{\log(2M/\delta')/(2n)}$ simultaneously --- in
particular for $\hat\pi$ and $\pi^\dagger$, wherever the walk
terminates. The deployed
$\hat\pi$ minimizes $\hat c$ over the certified set, which contains
$\pi^\dagger$, hence
$c(\hat\pi) \le \hat c(\hat\pi) + \eps_n \le \hat c(\pi^\dagger) +
\eps_n \le c(\pi^\dagger) + 2\eps_n$. Both events hold together with
probability $\ge 1 - 2\delta'$. (If the cheapest-selection is
restricted to the first $K$ walk positions, or costs are estimated on
an independent split, the union bound over $M$ shrinks to $K$ and
$\eps_n$ improves accordingly; validity is unaffected either way.)
Without the monotone-order premise the
walk may stop before $K$; the certified-prefix optimality and the
$2\eps_n$ cost bound continue to hold with $\pi^\dagger$ replaced by
the cheapest feasible candidate of the certified prefix, and validity
is unaffected. \hfill$\square$

\subsection{Proof of Proposition~\ref{prop:lifetime}}
Index epochs $k = 0, 1, 2, \dots$ in order of their terminating events
--- an alarm or the exhaustion of the monitor's horizon $T$ (scheduled
restart). Each epoch lasts at most $T$ queries, so an unbounded
lifetime is covered by countably many epochs, and each epoch's monitor
is a fresh finite-horizon detector run strictly within its own block,
the regime in which Proposition~\ref{prop:eprocess} holds. Epoch $k$'s
certification is a fixed-sequence test at level
$\delta_k/2$ run on data from that epoch alone --- the initial
certification is epoch 0's, run at $\delta_0/2 = \delta/4$ on the
initial calibration split, so it is inside the schedule rather than an
extra $\delta$-sized event on top of it --- and by
Theorem~\ref{thm:validity} the probability that epoch $k$ deploys a
policy violating its contract (on the epoch's distribution) is at most
$\delta_k/2$; epoch $k$'s monitor is an e-detector at level
$\delta_k/2$, so by Proposition~\ref{prop:eprocess} the probability
that it false-alarms while the epoch is compliant is at most
$\delta_k/2$. Horizon-triggered boundaries consume no additional error
probability: exhausting $T$ without alarm is a deterministic event,
and the subsequent recertification is already accounted at level
$\delta_{k+1}/2$. The lifetime error event is the union over $k$ of
these
$2$ events per epoch; by countable subadditivity its probability is at
most $\sum_{k\ge 0} \delta_k = \sum_{k\ge 0} \delta/2^{k+1} = \delta$,
with no bound on the number of epochs. For the delay claim, the
$k$-th monitor's alarm threshold is $\log(2^{k+2}/\delta) =
\log(1/\delta) + (k+2)\log 2$, which enters the bound of
Proposition~\ref{prop:delay} additively, giving the stated
$(k+2)\log 2 / g^\ast$ increment; a change straddling a block boundary
is caught by the next block's monitor, adding at most the current
block's residual to the delay. Under the infinite-horizon alternative
--- prior weights $\frac{1}{r(r+1)}$ over restart times, summing to
one --- validity needs no horizon, but the mixture mass at restart
$\tau$ is $\Theta(\tau^{-2})$, so the threshold term becomes
$\log(\tau^2 |W|)$ and the delay inherits a $\log\tau$ dependence;
the block construction avoids this. \hfill$\square$

\begin{rmk}[Audit subsampling]
Let each query be audited independently with probability $r$ and let
the monitor be run on the audited subsequence. Conditional on being
audited, a query is an unbiased draw from the deployment distribution
(the audit coin is independent of the query), so the audited losses
$(\ell_{t_j})_j$ satisfy the same conditional-mean hypotheses as the
full stream, and Propositions~\ref{prop:eprocess}--\ref{prop:delay}
apply verbatim with $t$ counting audited queries. A post-change window
of $N$ audited queries occupies $N/r$ wall-clock queries in
expectation (and concentrates at rate $O(\sqrt{N}/r)$), giving the
$1/r$ delay dilation quoted in Section~\ref{sec:monitor}.
\end{rmk}
